\chapter{Resultados de Complexidade em Grafos planares}

\alv{Alguma coisa aqui para introduzir}



\section{NP-completude para grafos planares orientados}

Em~\cite{CENTENO200911}, foi provado que no caso não orientado, decidir se o número de intervalo (\inum{P_3}) é no máximo $k$ é um problema NP-completo. Posteriormente, \cite{PENSO201583} mostraram que o problema continua NP-completo mesmo em grafos planares com grau máximo 3. Além disso, em\cite{PENSO201583}, é apresentado o problema \sam.

\vspace{5pt}

\problemDef{\sam}{Um conjunto de cláusulas $F = \{C_1, C_2, \ldots, C_m\}$, um conjunto finito de variáveis booleanas $X = \{x_1, x_2, \ldots, x_n\}$, onde cada cláusula contém no máximo 3 literais, cada variável aparece no máximo 3 vezes, cada literal ocorre no máximo 2 vezes e o grafo bipartido $H_F = (V, E)$ tal que $V = \{w_{c_1}, w_{c_2}, \ldots, w_{c_m}\} \cup \{v_{x_1}, v_{x_2}, \ldots, v_{x_n}\}$ e $E$ contém exatamente os pares $(w_{c_i}, v_{x_j})$ tal que $x_j$ ou $\overline{x_j}$ pertencem a cláusula $C_i$, é planar}{Existe alguma valoração de valores verdadeiros as variáveis de $X$ que satisfaz $F$?}{}{psam3}

\begin{theorem}[\cite{PENSO201583}]
    \sam é NP-completo
\end{theorem}

Muitos conceitos de convexidade em grafos se estendem aos grafos orientados. Tendo como base o resultado de \cite{PENSO201583} faremos uma redução do problema \sam para o problema do \ind.

\problemDef{\ind}{Um grafo planar orientado D e um inteiro positivo \emph{k}}{\inumd{P_3} $\leq k$?}{}{gpinp3}

\begin{theorem}
    \ind é \npc.
\begin{proof}
    Essa demonstração é similar ao resultado não direcionado apresentado em \cite{PENSO201583}. Podemos ver que \ind está em \np tendo um subconjunto $S$ de vértices de um grafo direcionado $D$ como certificado. Neste caso, basta verificar o tamanho do subconjunto $S$ e para cada $v \in V(D) \backslash S$, verificar se existe pelo menos um vértice $v_1 \in N^+(v)$ e pelo menos um vértice $v_2 \in N^-(v)$ que estão em $S$.

    Para provar a \npct, faremos uma redução do problema \sam para \ind. Seja $F$ uma instância de \sam, construíremos um grafo orientado $D$ da seguinte forma:

    \begin{itemize}
        \item Para cada variável $x_i$, criamos em $D$ um \emph{gadget} $d_{x_i}$ composta pelo ciclo direcionado $f^1_{x_i}, t^1_{x_i}, a^1_{x_i}, a^2_{x_i},f^2_{x_i}, t^2_{x_i}, a^3_{x_i}, a^4_{x_i}$;

        \item Para cada cláusula com no máximo 2 literais, criamos em $D$ um \emph{gadget} $d_{c_j}$ composta pelos vértices $c^1_j$ e $c^2_j$ e o arco $(c^2_j, c^1_j)$;
        
        \item Para cada cláusula com exatamente 3 literais, criamos em $D$ um \emph{gadget} $d_{x_i}$ composta pelos vértices $c^1_j, c^2_j, c^3_j, l^1_j, l^2_j$ e os arcos ($c^2_j, c^1_j$), ($c^1_j, c^3_j$), ($c^3_j, l^2_j$) e ($c^1_j, l^1_j$)

    \end{itemize}

    Considerando a representação plana do grafo $H_F$ como um \emph{template} para o grafo orientado $D$, adicionamos arcos que conectam os \emph{gadgets} da seguinte forma: 

    \begin{itemize}
        \item Se a variável $x_i$ aparece positiva (respectivamente, negativa) na cláusula $C_j$, adicionamos o arco ($t^a_{x_i}, c^b_j$) (respectivamente, ($f^a_{x_i}, c^b_j$)), onde $a \in \{1,2\}$ e $b \in \{2,3\}$.
        \item Se a cláusula contém no máximo 2 literais, o único vértice conector é o vértice $c^2_j$. Por outro lado, se a cláusula tem exatamente 3 literais, os conectores são apenas $c^2_j$ e $c^3_j$. Neste caso, $c^2_j$ terá 2 vizinhos de entrada e $c^3_j$ terá exatamente 1 vizinho de entrada.
        
        \alv{a única forma de garantir isso é tendo grau de entrada no máximo 2?}

        \item Neste caso, o grafo orientado $D$ tem grau de entrada máximo igual a 2
    \end{itemize}

Vamos demonstrar que dada uma instância $F$ de \sam, onde $n$ é o número de variáveis, $m_1$ é o número de cláusulas com no máximo 2 literais e $m_2$ é o número de cláusulas com exatamente 3 literais, $F$ é satisfatível se e somente se o grafo orientado $D$ obtido na construção acima tem um conjunto de intervalo \pd de tamanho $k = 4n + m_1 + 3m_2$.

Seja $F$ uma fórmula \sam e $A$ uma atribuição de valores verdadeiros para $F$. Construímos um conjunto de intervalo $S_A$ da seguinte forma:

\begin{enumerate}
    \item Todos os vértices extremos são adicionados a $S_A$
    \item Se $x_i = verdadeiro$ em $A$, adicionamos os seguintes vértices a $S_A$: $t^1_{x_i}, t^2_{x_i}, a^2_{x_i}, a^4_{x_i}$
    \item Se $x_i = falso$ em $A$, adicionamos os seguintes vértices a $S_A$: $f^1_{x_i}, f^2_{x_i}, a^1_{x_i}, a^3_{x_i}$
    \item Para cada cláusula com exatamente 3 literais, seja $v \neq c^1_j$, se o vértice $v \in N^-(c^3_j)$ está em $S_A$ , então adicionamos o vértice $c^2_j$ a $S_A$, caso contrário, adicionamos o vértice $c^1_j$.
    \item Como $A$ é uma atribuição de valores verdade, cada cláusula deve ser satisfeita, com isso temos que cada \emph{gadget} $d_{c_j}$ de $D$ deve ter pelo menos um vizinho no conjunto de intervalo $S_A$ que estão em algum $d_{x_i}$. Portanto, $S_A$ é um conjunto de intervalo tal que $|S_A| = 4n+m_1+3m_2$.
\end{enumerate}

Por outro lado, seja $S_A \subseteq V(D)$ um conjunto de intervalo tal que $|S_A| = 4n+m_1+3m_2 = k$, obtemos uma atribuição de valores $A$ da seguinte forma:
\begin{enumerate}
    \item Todo vértice extremo está em $S_A$ \alv{usar a referencia para garantir a escrita dessa forma}
    \item Pelo menos um vértice do \emph{gadget} $d_{c_j}$ quando a cláusula tem até 2 literais deve estar em $S_A$ e pelo menos 3 vértices do \emph{gadget} $d_{c_j}$ quando a cláusula tem exatamente 3 literais devem estar em $S_A$. Como $|S_A| = k$, temos que exatamente 4 vértices do \emph{gadget} $d_{x_i}$ devem estar em $S_A$. Pela construção do digrafo $D$, vemos que as únicas configurações possíveis são: Adicionar os vértices $\{t^1_{x_i},t^2_{x_i},a^2_{x_i},a^4_{x_i}\}$ ao conjunto $S_A$ ou adicionar os vértices $\{f^1_{x_i},f^2_{x_i},a^1_{x_i},a^3_{x_i}\}$.
    \item Com isso, $x_i = verdadeiro$ se e somente se $\{t^1_{x_i},t^2_{x_i},a^2_{x_i},a^4_{x_i}\} \in S_A$. Desta forma, temos que cada \emph{gadget} $d_{c_j}$ deve ter um vizinho em $S_A$, se não, $S_A$ não seria um conjunto de intervalo do digrafo $D$ ou $|S_A| > k$, já que mais vértices em cada \emph{gadget} de cláusula estariam em $S_A$. Portanto, se $D$ tem um conjunto de intervalo de tamanho $4n+m_1+3m_2$, então, a fórmula $F$ é satisfeita.
\end{enumerate}

\end{proof}
\end{theorem}

Naturalmente, temos que além do número de intervalo, outro parâmetro bastante estudado é o número de envoltória que também se estende aos grafos orientados. Tendo com base o resultado de \cite{PENSO201583} faremos uma redução do problema \sam para o problema \hnd.

\problemDef{\hnd}{Um grafo planar orientado D e um inteiro positivo \emph{k}}{\hnumd{P_3} $\leq k$?}{}{gphnp3}

\begin{theorem}
    \hnd é \npc
    \begin{proof}
        
    \end{proof}
\end{theorem}

